%==========================================================
%  CHAPTER 1: The Canadian Economy — Structure, Institutions,
%              and Analytical Tools
%  ECON 204 / Contemporary Economic Issues: A Canadian Perspective
%  Muhebullah Karimzada — University of Northern British Columbia
%==========================================================

\chapter{The Canadian Economy: Structure, Institutions, and Analytical Tools}
\label{ch:introduction}

%----------------------------------------------------------
%  OPENING VIGNETTE
%----------------------------------------------------------

\noindent In the spring of 2024, four people in Prince George, British Columbia
faced economic decisions shaped by forces none of them fully understood.

Diane, 52, received her mortgage renewal notice in March. Her monthly payment
was rising from \$1,840 to \$2,710 --- an increase of \$870, or more than
\$10,000 per year. Nothing had changed in her household except one number:
the Bank of Canada's overnight interest rate, which had climbed from 0.25\%
in early 2022 to 5.00\% by mid-2023. She had not overspent. She had not lost
her job. But a decision made in an Ottawa boardroom had fundamentally altered
her monthly budget.

Her son Marcus, 24, had just graduated from UNBC with a degree in business
administration. He was earning \$22 per hour at a local firm --- a wage that
would have seemed generous five years ago, but which, after two years of
inflation averaging over 5\%, felt significantly less secure. He was
trying to save for a down payment on an apartment, but rental costs in Prince
George had risen nearly 30\% since 2020. His parents' house had doubled in
value, but that wealth would not be his for decades.

At the same sawmill where Marcus's uncle worked, the plant manager was
navigating a \$3 million investment decision: automate the log-grading line
with computer vision technology, or keep 18 workers. The system would pay for
itself in four years. The workers earned above the minimum wage, had families
in the community, and spent their incomes at local businesses. The economics
pointed one direction; the community consequences pointed another.

And at the Northern Health Authority, administrator Priya was reviewing a
budget that had grown 14\% in two years --- driven by an aging patient
population, nurse recruitment costs, and pharmaceutical price increases ---
while the provincial government's transfer payment had grown by only 6\%.
The gap had to come from somewhere.

These four people live in the same economy. They are affected by the same
interest rates, the same inflation, the same labour market pressures, and the
same fiscal constraints. None of their situations can be understood without
understanding how the Canadian economy works, how its institutions make
decisions, and how those decisions ripple outward into ordinary lives.

That is what this textbook is about.

%----------------------------------------------------------
%  LEARNING OBJECTIVES
%----------------------------------------------------------

\begin{learningobjectives}
  \item Describe the scale and structure of the Canadian economy using key
        aggregate statistics, including GDP, GDP per capita, and sectoral
        composition.
  \item Apply the expenditure approach to GDP and calculate nominal GDP growth,
        real GDP growth, and GDP per capita from raw data.
  \item Identify Canada's major economic institutions and explain the mandate,
        tools, and independence of each.
  \item Distinguish between monetary policy and fiscal policy, explain the
        transmission channels of each, and identify which institution is
        responsible for each in Canada.
  \item Apply the core tools of economic reasoning --- models, opportunity cost,
        marginal analysis, and the positive--normative distinction --- to
        real-world Canadian policy questions.
  \item Situate Canada within an international context by comparing key
        macroeconomic indicators with OECD peers.
  \item Preview the ten major contemporary economic issues addressed in this
        textbook and explain how they are interconnected.
\end{learningobjectives}

%==========================================================
\section{Why Study Contemporary Economic Issues?}
%==========================================================

Economics is often introduced as the study of scarcity: because resources are
limited and human wants are not, societies must make choices about how to
allocate what they have. That definition is correct, but it understates the
subject's reach. A more precise framing for this course is that economics is
the study of \term{incentives}, \term{trade-offs}, and \term{unintended
consequences} --- the three forces that shape virtually every economic outcome
a society produces.

\term{Contemporary economic issues} are economic problems that are actively
unfolding, contested among reasonable people, and consequential for the
well-being of Canadians. They differ from purely theoretical questions
because they involve real institutions making real decisions under genuine
uncertainty. They differ from purely political questions because the
empirical evidence matters --- how markets actually behave, how people
actually respond to incentives, and what policies have actually worked
in comparable contexts.

Consider the four people from our opening vignette. Diane's mortgage crisis
cannot be understood without knowing why the Bank of Canada raised interest
rates, what caused inflation to spike in 2022, and how monetary policy
transmits through the financial system to household balance sheets. Marcus's
wage frustration cannot be understood without knowing how labour markets
determine compensation, why real wages have lagged productivity growth since
the 1980s, and how automation changes the demand for different kinds of work.
The sawmill manager's dilemma cannot be understood without the economics of
technological adoption, capital investment, and labour market adjustment. And
Priya's budget cannot be balanced without understanding fiscal federalism,
demographic aging, and the political economy of healthcare spending.

Economics provides the analytical tools to work through each of these problems
systematically. It does not always deliver clean, comfortable answers. But it
imposes a discipline on reasoning that is absent from purely ideological or
intuitive approaches: it forces you to trace consequences, identify trade-offs,
acknowledge trade-offs on the other side, and ask whether the evidence
supports your conclusions.

\begin{thinkbox}
When you encounter a policy proposal in this course --- a carbon tax, a rent
control ordinance, an immigration target, a minimum wage increase --- train
yourself to ask three questions before forming an opinion:
\begin{enumerate}[(1)]
  \item \textbf{What incentives does this policy create?} Who responds, how,
        and in what direction?
  \item \textbf{What are the trade-offs?} What must be given up to achieve
        this goal?
  \item \textbf{What might happen that the policymaker did not intend?}
        Economics is full of well-intentioned policies that produced
        surprising consequences.
\end{enumerate}
These questions do not determine whether a policy is good or bad. They are
the beginning of a serious analysis, not its conclusion.
\end{thinkbox}

%==========================================================
\section{The Canadian Economy at a Glance}
%==========================================================

Before analysing specific contemporary issues, we need a clear statistical
portrait of the economy we are studying. This section establishes the
foundational facts of Canadian economic structure that every subsequent
chapter will draw upon.

%----------------------------------------------------------
\subsection{Scale and Aggregate Output}
%----------------------------------------------------------

Canada is the world's ninth-largest economy measured by nominal gross domestic
product. In 2024, Canada's GDP reached approximately \$3.0 trillion in
Canadian dollars --- a figure that places Canada ahead of Australia, South
Korea, and Spain, but significantly behind the United States (\$29~trillion
USD) and far smaller than China and the European Union collectively. With a
population of approximately 41.5 million people, Canada's GDP per capita ---
our first approximation of average living standards --- was roughly
\$72,000~CAD, or approximately \$53,000~USD at prevailing exchange rates.

\begin{defbox}{Gross Domestic Product (GDP)}
\textbf{GDP} is the total market value of all final goods and services
produced within a country's borders during a given period, typically one
calendar year or one fiscal quarter.

\textit{Economic intuition:} GDP measures the size of the economy's
productive output. It captures what the economy \textit{produces} --- not
what it owns, not what it earns abroad, and not what flows from one
person to another without new production (such as transfer payments or
the sale of existing assets).

\textit{What GDP excludes:} Household production (cooking, childcare at
home), the underground economy, environmental degradation, the distribution
of income, and non-material well-being.

\textit{Canadian example:} In 2020, Canada's real GDP fell 5.2\% ---
the steepest single-year contraction since the Great Depression --- as
COVID-19 forced businesses to close and household spending collapsed.
By 2021, GDP rebounded by 5.0\% as restrictions lifted. The swing of
more than 10 percentage points in two years illustrates both the fragility
and the resilience of a modern market economy.
\end{defbox}

%----------------------------------------------------------
\subsection{The Expenditure Approach to GDP}
\label{sec:expenditure}
%----------------------------------------------------------

There are three equivalent ways to measure GDP: the expenditure approach,
the income approach, and the production approach. In practice, national
statistical agencies use all three and reconcile any discrepancies.
For macroeconomic analysis, the \term{expenditure approach} is the most
widely used because it decomposes GDP by who is spending.

The fundamental expenditure identity is:

\begin{equation}
  Y \;=\; C \;+\; I \;+\; G \;+\; (X - M)
  \label{eq:gdp_expenditure}
\end{equation}

\noindent where:
\begin{itemize}
  \item $Y$ is nominal GDP (total output, measured in current dollars)
  \item $C$ is \term{household consumption expenditure} --- spending by
        households on goods (food, clothing, vehicles) and services
        (healthcare, education, financial services). This is typically
        the largest component, accounting for roughly 55--60\% of GDP
        in Canada.
  \item $I$ is \term{gross investment} (or gross capital formation) ---
        business spending on plant, equipment, and inventory, plus
        residential construction. This is the most volatile component
        of GDP.
  \item $G$ is \term{government expenditure on goods and services} ---
        note this includes only direct government purchases, \textit{not}
        transfer payments such as Employment Insurance or the Canada
        Child Benefit, which are redistribution rather than new production.
  \item $X$ is exports --- goods and services produced in Canada and
        sold abroad.
  \item $M$ is imports --- goods and services produced abroad and
        purchased in Canada.
  \item $(X - M) = NX$ is \term{net exports} (the trade balance).
        A positive NX means Canada sells more abroad than it buys ---
        a trade surplus. A negative NX means a trade deficit.
\end{itemize}

\begin{databox}{Canada's GDP by Expenditure Component, 2024 (Approximate)}
\begin{center}
\renewcommand{\arraystretch}{1.4}
\begin{tabular}{lrr}
\toprule
\textbf{Component} & \textbf{Value (\$B~CAD)} & \textbf{Share of GDP (\%)} \\
\midrule
Household Consumption ($C$)     & 1{,}652 & 55.1 \\
Gross Investment ($I$)          &   652  & 21.7 \\
Government Expenditure ($G$)    &   718  & 23.9 \\
Net Exports ($NX = X - M$)      &  $-$21 & $-$0.7 \\
\midrule
\textbf{GDP} ($Y = C+I+G+NX$)   & \textbf{3{,}001} & \textbf{100.0} \\
\bottomrule
\end{tabular}
\end{center}
\vspace{2pt}
\textit{Note:} Values are approximate and illustrative, consistent with
Statistics Canada national accounts data. Figures may not sum to exactly
3{,}001 due to rounding.
\source{Statistics Canada, Table 36-10-0104-01 (National Accounts, expenditure
approach), 2024 estimates.}
\end{databox}

\noindent\textbf{Worked Example 1.1 --- Calculating GDP from Components}

Suppose you are given the following data for a simplified economy:

\begin{center}
\begin{tabular}{lr}
\toprule
Item & Value (\$B) \\
\midrule
Household spending on food, housing, and services & 1{,}200 \\
Business investment in machinery and factories   &   400  \\
Residential housing construction                  &   150  \\
Federal, provincial, and municipal government purchases & 500 \\
Exports of goods and services                    &   620  \\
Imports of goods and services                    &   670  \\
\bottomrule
\end{tabular}
\end{center}

\noindent\textit{Step 1:} Identify the components of Equation~\eqref{eq:gdp_expenditure}:
\begin{align*}
C &= 1{,}200 \\
I &= 400 + 150 = 550 \quad \text{(business investment + residential construction)} \\
G &= 500 \\
NX &= 620 - 670 = -50 \quad \text{(a trade deficit of \$50 billion)}
\end{align*}

\noindent\textit{Step 2:} Apply the identity:
\[
Y = 1{,}200 + 550 + 500 + (-50) = \mathbf{2{,}200 \text{ billion}}
\]

\noindent\textit{Interpretation:} This economy produces \$2.2~trillion of
output. The trade deficit (\$50B) means the economy is spending more than
it produces --- the excess is financed by borrowing from abroad or by
foreigners purchasing domestic assets.

\begin{keytakeaway}
GDP measures production, not welfare. A society can have high GDP and
high inequality, high GDP and poor health outcomes, or high GDP and
severe environmental degradation. GDP is an indispensable starting point
for macroeconomic analysis, but it is not the endpoint for evaluating
economic well-being.
\end{keytakeaway}

%----------------------------------------------------------
\subsection{Real GDP, Nominal GDP, and the GDP Deflator}
\label{sec:realgdp}
%----------------------------------------------------------

When we compare GDP across years, we face a measurement challenge: prices
change over time. If Canadian GDP rose from \$2.8~trillion in 2022 to
\$3.0~trillion in 2024, some of that increase reflects more goods and
services being produced, but some reflects higher prices for the same
goods and services. Only the first part represents genuine economic growth.

\term{Nominal GDP} measures output at the \textit{current year's prices}.
\term{Real GDP} measures output at the prices of a fixed \term{base year},
removing the effect of inflation.

The conversion between them uses the \term{GDP deflator}:

\begin{equation}
  \text{GDP Deflator}_t = \frac{\text{Nominal GDP}_t}{\text{Real GDP}_t} \times 100
  \label{eq:gdp_deflator}
\end{equation}

Rearranging, real GDP can be obtained from nominal GDP and the deflator:

\begin{equation}
  \text{Real GDP}_t = \frac{\text{Nominal GDP}_t}{\text{GDP Deflator}_t / 100}
  \label{eq:real_gdp}
\end{equation}

The \term{real GDP growth rate} --- the standard measure of economic
expansion or contraction --- is:

\begin{equation}
  g_t = \frac{\text{Real GDP}_t - \text{Real GDP}_{t-1}}{\text{Real GDP}_{t-1}} \times 100
  \label{eq:gdp_growth}
\end{equation}

\noindent\textbf{Worked Example 1.2 --- Real vs. Nominal GDP Growth, Canada 2022--2023}

\begin{center}
\begin{tabular}{lrrrr}
\toprule
Year & Nominal GDP (\$B) & GDP Deflator & Real GDP (\$B, 2017\$) & Real Growth (\%) \\
\midrule
2022 & 2{,}721 & 110.3 & 2{,}467 & 3.4 \\
2023 & 2{,}879 & 115.2 & 2{,}499 & 1.3 \\
\bottomrule
\end{tabular}
\end{center}

\noindent\textit{Calculation for 2023 real GDP:}
\[
\text{Real GDP}_{2023} = \frac{2{,}879}{115.2/100} = \frac{2{,}879}{1.152} \approx 2{,}499 \text{ billion (2017 dollars)}
\]

\noindent\textit{Calculation for real growth rate:}
\[
g_{2023} = \frac{2{,}499 - 2{,}467}{2{,}467} \times 100 = \frac{32}{2{,}467} \times 100 \approx 1.3\%
\]

\noindent\textit{Interpretation:} Nominal GDP rose 5.8\% between 2022
and 2023, but more than four percentage points of that increase reflected
higher prices rather than more output. Real growth was only 1.3\% ---
an economy moving at a moderate pace, consistent with Canada's estimated
potential growth rate of approximately 1.5--2.0\% per year in the
mid-2020s.

%----------------------------------------------------------
\subsection{GDP Per Capita: A Measure of Living Standards}
%----------------------------------------------------------

Total GDP tells us the size of the economy; it tells us very little about
the typical person's standard of living. An economy of 1.4 billion people
will have a larger GDP than one of 40 million people, even if each of its
citizens is far poorer. For cross-country and cross-time comparisons
of living standards, economists use \term{GDP per capita} --- total output
divided by population:

\begin{equation}
  \text{GDP per capita}_t = \frac{Y_t}{N_t}
  \label{eq:gdp_per_capita}
\end{equation}

\noindent where $N_t$ is the total population.

\noindent\textbf{Worked Example 1.3 --- Canada vs. United States GDP per Capita, 2024}

\begin{center}
\renewcommand{\arraystretch}{1.3}
\begin{tabular}{lrrr}
\toprule
Country & Nominal GDP (USD B) & Population (M) & GDP per Capita (USD) \\
\midrule
United States & 28{,}781 & 337.0 & 85{,}406 \\
Canada        & 2{,}221  &  41.5 & 53{,}518 \\
\midrule
\textbf{Canada--US Gap} & & & \textbf{--\$31{,}888 (--37.3\%)} \\
\bottomrule
\end{tabular}
\end{center}

\noindent\textit{Note:} Canada's nominal GDP is converted to USD at
approximate 2024 exchange rate of 0.74 USD/CAD. GDP in USD thousands:
\(\$2{,}221\text{B} \times 0.74 = \$1{,}643\text{B}\)\ USD; \(\$1{,}643\text{B} / 41.5\text{M} = \$39{,}590\) USD. For PPP-adjusted comparison (which removes exchange rate distortions), the gap narrows but Canada still trails the US by approximately 20--25\%.

This income gap with the United States is one of the most important and
persistent facts about the Canadian economy. It is not primarily explained
by natural resources, education, or geography. A major contributor is
a productivity gap --- Canadians produce less output per hour worked than
Americans --- and that productivity puzzle is the subject of
Chapter~\ref{ch:technology}.

\begin{figure}[H]
\centering
\begin{tikzpicture}
\begin{axis}[
  ybar,
  bar width=0.55cm,
  width=13cm, height=7.5cm,
  symbolic x coords={Italy, Japan, {OECD Avg}, France, UK, Canada,
                     Germany, Australia, Switzerland, USA},
  xtick=data,
  x tick label style={rotate=35, anchor=east, font=\small},
  ylabel={GDP per Capita, PPP-Adjusted (USD thousands), 2024},
  ylabel style={font=\small},
  ymin=0, ymax=100,
  ytick={0,20,40,60,80,100},
  yticklabel style={font=\small},
  ymajorgrids=true,
  grid style={dashed, gray!25},
  enlarge x limits=0.08,
  nodes near coords,
  nodes near coords style={font=\footnotesize, rotate=90, anchor=west},
  every axis plot/.append style={fill=unbcgreen!70, draw=unbcgreen!90}
]
\addplot coordinates {
  (Italy,        46.0)
  (Japan,        52.0)
  ({OECD Avg},   52.5)
  (France,       56.5)
  (UK,           55.0)
  (Canada,       57.0)
  (Germany,      62.0)
  (Australia,    64.5)
  (Switzerland,  78.0)
  (USA,          85.4)
};
\end{axis}
\end{tikzpicture}
\caption{GDP per Capita (PPP-Adjusted, USD thousands), Selected OECD Countries, 2024.
Canada's living standard exceeds the OECD average and surpasses France, the UK,
Japan, and Italy, but trails Germany, Australia, Switzerland, and the United
States. The large Canada--US gap of approximately 37\%, despite similar geography,
institutions, and education levels, reflects a structural productivity difference
explored in Chapter~\ref{ch:technology}.}
\label{fig:gdp_per_capita}
\source{OECD National Accounts at a Glance, 2024; World Bank World Development Indicators.}
\end{figure}

%----------------------------------------------------------
\subsection{Sectoral Structure: The Service Economy}
%----------------------------------------------------------

Canada's economy is predominantly a \textbf{service economy}. Services ---
finance, real estate, retail, healthcare, professional services, education,
and public administration --- account for approximately 77\% of GDP. This
pattern is common among high-income countries: as economies develop, labour
and capital shift from agriculture to manufacturing to services.

\begin{figure}[H]
\centering
\begin{tikzpicture}
\begin{axis}[
  ybar,
  bar width=0.50cm,
  width=14cm, height=7cm,
  symbolic x coords={{Retail}, {Public Admin}, {Healthcare}, {Construction},
                     {Transport}, {Prof.\ Svcs}, {Mining/Oil}, {Manufacturing},
                     {Finance}, {Real Estate}},
  xtick=data,
  x tick label style={rotate=38, anchor=east, font=\small},
  ylabel={Share of GDP (\%)},
  ylabel style={font=\small},
  ymin=0, ymax=16,
  ytick={0,2,4,6,8,10,12,14,16},
  yticklabel style={font=\small},
  ymajorgrids=true,
  grid style={dashed, gray!25},
  enlarge x limits=0.06,
  every axis plot/.append style={fill=unbcgreen!65, draw=unbcgreen!85}
]
\addplot coordinates {
  ({Retail},       5.5)
  ({Public Admin}, 6.0)
  ({Healthcare},   6.2)
  ({Construction}, 7.3)
  ({Transport},    5.4)
  ({Prof.\ Svcs},  7.0)
  ({Mining/Oil},   9.1)
  ({Manufacturing},10.2)
  ({Finance},      7.8)
  ({Real Estate},  13.5)
};
\end{axis}
\end{tikzpicture}
\caption{Canada's Major Industries as a Share of GDP, 2024.
Real estate dominates --- reflecting high property values and the
construction and rental sectors --- while manufacturing and mining/oil
and gas follow. Mining and oil and gas (9.1\%) punches above its weight
politically relative to its actual GDP share, a tension explored in
Chapter~\ref{ch:climate}.
The service-heavy structure means Canada's business cycle is largely
driven by consumer spending, housing, and financial conditions --- all
channels through which monetary policy operates powerfully.}
\label{fig:gdp_by_industry}
\source{Statistics Canada, Table 36-10-0401-01, GDP at basic prices by industry, 2024.}
\end{figure}

Critically, within the goods-producing sector, \textbf{natural resources}
(oil and gas, mining, forestry, and agriculture) play a disproportionately
large role relative to most comparable advanced economies. Canada's oil sands
in Alberta produce over 3 million barrels of crude oil per day. British
Columbia's forest products sector exports billions of dollars of softwood lumber
annually. Saskatchewan supplies approximately 40\% of the world's potash.
These resources generate export revenues, government royalties, and employment ---
but they also make Canada's economy sensitive to global commodity price cycles
and they are central to the climate transition challenge explored in
Chapter~\ref{ch:climate}.

%----------------------------------------------------------
\subsection{Regional Economic Geography}
%----------------------------------------------------------

Canada is a vast country, and its economy is not evenly distributed. The
regional concentration of economic activity creates sharp differences in
employment conditions, wages, housing markets, and fiscal capacity across
provinces --- differences that make national averages sometimes misleading
for understanding local economic conditions.

\begin{figure}[H]
\centering
\begin{tikzpicture}
\begin{axis}[
  ybar,
  bar width=0.45cm,
  width=14cm, height=7cm,
  symbolic x coords={{Terr.}, {PEI}, {N\!L}, {NB}, {NS},
                     {MB}, {SK}, {BC}, {AB}, {QC}, {ON}},
  xtick=data,
  x tick label style={font=\small},
  ylabel={Share of National GDP (\%)},
  ylabel style={font=\small},
  ymin=0, ymax=44,
  ytick={0,5,10,15,20,25,30,35,40},
  yticklabel style={font=\small},
  ymajorgrids=true,
  grid style={dashed, gray!25},
  enlarge x limits=0.06,
  nodes near coords,
  nodes near coords style={font=\footnotesize, rotate=90, anchor=west},
  every axis plot/.append style={fill=policyblue!60, draw=policyblue!80}
]
\addplot coordinates {
  ({Terr.}, 0.5)
  ({PEI},   0.4)
  ({N\!L},  1.2)
  ({NB},    1.4)
  ({NS},    1.6)
  ({MB},    3.0)
  ({SK},    3.1)
  ({BC},   13.6)
  ({AB},   15.9)
  ({QC},   21.1)
  ({ON},   38.9)
};
\end{axis}
\end{tikzpicture}
\caption{Share of National GDP by Province and Territory, 2024.
Ontario and Quebec together account for 60\% of Canadian economic output.
Alberta, despite having one-quarter of Ontario's population, generates 15.9\%
of GDP --- reflecting the outsized economic weight of the oil sands. The three
territories combined produce 0.5\% of GDP despite covering 40\% of Canada's
land mass. This extreme geographic concentration of economic activity motivates
the equalization payment system examined in Chapter~\ref{ch:fiscal}.}
\label{fig:gdp_by_province}
\source{Statistics Canada, Provincial and Territorial Economic Accounts, 2024.}
\end{figure}

The regional diversity of Canada's economy creates significant policy
challenges. When oil prices fall, Alberta's provincial revenue collapses
--- but Ontario's economy is largely unaffected. When the Bank of Canada
raises interest rates to combat national inflation, the impact on heavily
indebted households in Toronto and Vancouver is far more severe than on
renters in Moncton or resource workers in Fort McMurray. National policy
instruments applied to a regionally diverse economy inevitably produce
unequal regional consequences.

\begin{keytakeaway}
Canada's economy is large in aggregate but regionally unequal: Ontario
and Quebec dominate output; Alberta dominates resources; and the Atlantic
provinces and territories depend heavily on federal transfers. National
economic statistics always need to be disaggregated by region for a
complete picture.
\end{keytakeaway}

%==========================================================
\section{Canada's Trade Dependence}
\label{sec:trade_dependence}
%==========================================================

Canada is one of the most trade-dependent large economies in the world.
\term{Trade openness} --- defined as total trade (exports plus imports) as
a share of GDP --- measures how integrated an economy is with the global
trading system:

\begin{equation}
  \text{Trade Openness} = \frac{X + M}{Y} \times 100
  \label{eq:trade_openness}
\end{equation}

For Canada in 2024:
\[
\text{Trade Openness} = \frac{897 + 918}{3{,}001} \times 100 = \frac{1{,}815}{3{,}001} \times 100 \approx 60.5\%
\]

This means that the sum of Canada's exports and imports is equivalent
to more than 60\% of total GDP --- an extraordinary degree of integration
for an economy of Canada's size. By comparison, the United States, with
a much larger domestic market, has a trade openness ratio of approximately
27\%.

\begin{figure}[H]
\centering
\begin{tikzpicture}
\begin{axis}[
  ybar,
  bar width=0.6cm,
  width=12cm, height=7cm,
  symbolic x coords={USA, Japan, Italy, France, UK, Canada, Germany},
  xtick=data,
  x tick label style={font=\small},
  ylabel={Trade Openness: $(X+M)/Y \times 100$ (\%)},
  ylabel style={font=\small},
  ymin=0, ymax=105,
  ytick={0,20,40,60,80,100},
  yticklabel style={font=\small},
  ymajorgrids=true,
  grid style={dashed, gray!25},
  enlarge x limits=0.10,
  nodes near coords,
  every node near coord/.append style={font=\footnotesize},
  every axis plot/.append style={fill=exampleteal!65, draw=exampleteal!85}
]
\addplot coordinates {
  (USA,     27.0)
  (Japan,   38.0)
  (Italy,   56.0)
  (France,  58.0)
  (UK,      62.0)
  (Canada,  60.5)
  (Germany, 88.0)
};
\end{axis}
\end{tikzpicture}
\caption{Trade Openness --- $(X+M)/Y \times 100$ --- G7 Countries, 2024.
Canada's trade openness (60.5\%) is among the highest in the G7, surpassed
only by Germany. The US is far less trade-dependent because its enormous
domestic market satisfies most of its own demand. Canada's high trade
dependence means it is especially vulnerable to US tariff policy, exchange
rate movements, and global commodity price shocks --- themes developed in
Chapter~\ref{ch:trade}.}
\label{fig:trade_openness}
\source{World Bank World Development Indicators; Statistics Canada, Table 12-10-0121-01.}
\end{figure}

The most critical structural fact about Canada's trade is its \textbf{extreme
concentration} in a single market. The United States absorbs approximately
75\% of Canada's merchandise exports, making Canada more dependent on one
trading partner than almost any other advanced economy in the world. This
concentration is a source of enormous prosperity --- geographic proximity,
a shared language and legal tradition, and preferential market access under
CUSMA lower transaction costs dramatically --- but it is also a source of
strategic vulnerability.

\begin{table}[H]
\centering
\caption{Canada's Top 10 Trading Partners, 2024}
\label{tab:trading_partners}
\renewcommand{\arraystretch}{1.35}
\begin{tabular}{lrrr}
\toprule
\textbf{Country} & \textbf{Export Share (\%)} & \textbf{Import Share (\%)}
                 & \textbf{Trade Balance (\$B~CAD)} \\
\midrule
United States    & 75.4 & 51.2 & +130.2 \\
China            &  4.2 & 12.8 & $-$37.9 \\
United Kingdom   &  2.8 &  2.1 & $+$3.1 \\
Japan            &  2.1 &  2.8 & $-$4.0 \\
Germany          &  0.8 &  2.8 & $-$7.1 \\
Mexico           &  0.8 &  4.4 & $-$13.8 \\
South Korea      &  0.8 &  1.5 & $-$2.3 \\
Netherlands      &  1.2 &  0.6 & $+$2.2 \\
India            &  0.7 &  1.0 & $-$1.1 \\
France           &  0.6 &  1.2 & $-$2.4 \\
\midrule
\textbf{All others} & 10.6 & 19.6 & $-$30.5 \\
\bottomrule
\end{tabular}
\source{Statistics Canada, Table 12-10-0121-01; Global Affairs Canada,
Trade Data Online, 2024.}
\end{table}

\noindent Table~\ref{tab:trading_partners} reveals two important asymmetries.
First, Canada runs a substantial trade surplus with the United States (\$130B),
driven almost entirely by energy exports and auto-sector trade. Second, Canada
runs trade deficits with most other major partners, including China and Germany.
The concentration of the surplus in one country and one set of sectors creates
the vulnerability analysed in Chapter~\ref{ch:trade}: a shift in US trade
policy --- such as the imposition of tariffs on Canadian steel, lumber, or
automobiles --- can rapidly transform Canada's trade balance from surplus
to deficit.

%==========================================================
\section{Canada's Major Economic Institutions}
\label{sec:institutions}
%==========================================================

Economic outcomes do not emerge spontaneously from markets alone. They are
shaped by institutions --- the formal rules, organizations, and governance
frameworks within which economic decisions are made. Understanding Canada's
key economic institutions is a prerequisite for understanding how economic
policy is made and why its effects are what they are.

%----------------------------------------------------------
\subsection{The Bank of Canada}
%----------------------------------------------------------

The \term{Bank of Canada} is Canada's central bank. Established under the
\textit{Bank of Canada Act} in 1934, it is structured as a Crown corporation
that operates at formal arm's length from the elected federal government.
In practice, this independence is one of the most consequential institutional
features of the Canadian economy.

The Bank of Canada's primary mandate is to \textbf{promote the economic and
financial welfare of Canada}. Since 1991, this mandate has been operationalized
through a formal \term{inflation-targeting framework}, renewed most recently
in 2021: the Bank commits to keeping CPI inflation within a \textbf{1--3\%
control range}, with a \textbf{2\% midpoint target}.

The Bank conducts monetary policy primarily through its \term{policy interest
rate} --- the target for the \term{overnight rate}, the interest rate at which
major financial institutions lend to each other on an overnight basis. When
the Bank raises the overnight rate, it becomes more expensive to borrow
throughout the financial system. Mortgage rates rise. Business loans cost more.
The exchange rate tends to appreciate. All of these effects reduce spending
and investment, which eventually reduces inflation. When the Bank lowers the
overnight rate, the opposite occurs.

\begin{canexample}{The 2022--2024 Rate-Hike Cycle: The Fastest in Modern History}
Between March 2022 and July 2023, the Bank of Canada raised its policy
interest rate from 0.25\% to 5.00\% --- an increase of 475 basis points
(4.75 percentage points) in just 16 months. This was the most aggressive
monetary tightening cycle in the Bank of Canada's modern history, undertaken
to combat inflation that had reached 8.1\% in June 2022, the highest reading
since September 1983.

The rate hike cycle illustrates the transmission mechanism of monetary policy
in practice. By late 2023, inflation had fallen to approximately 3\%. By mid-2024,
it had approached the 2\% target. But the rate hikes also raised mortgage
payments for millions of Canadian households --- an estimated \$700 on average
for variable-rate mortgage holders --- and contributed to a sharp slowdown in
housing market activity. The Bank began cutting rates in June 2024.

This episode is the central case study of Chapter~\ref{ch:inflation}. It also
directly shapes the housing analysis in Chapter~\ref{ch:housing} and the
distributional analysis in Chapter~\ref{ch:inequality}.
\end{canexample}

%----------------------------------------------------------
\subsection{Department of Finance Canada}
%----------------------------------------------------------

The \term{Department of Finance Canada} is the federal government ministry
responsible for \term{fiscal policy} --- the use of government spending and
taxation to influence the economy. The Minister of Finance presents an
annual federal budget, typically in March, which outlines revenue plans,
spending priorities, and projections for the federal deficit or surplus.

Unlike the Bank of Canada, the Department of Finance operates under direct
democratic accountability: its decisions are made by elected politicians and
must pass through Parliament. This creates a fundamentally different set of
incentives than those governing the Bank of Canada, and it is one reason why
monetary and fiscal policy are deliberately separated in Canada's institutional
design.

Key agencies operating alongside Finance Canada include:

\begin{itemize}
  \item \textbf{Parliamentary Budget Officer (PBO):} An independent officer
        of Parliament established in 2006, mandated to provide non-partisan
        economic and fiscal analysis. The PBO has become an important check
        on government budget projections, regularly producing independent
        assessments that sometimes differ substantially from government
        forecasts.
  \item \textbf{Canada Revenue Agency (CRA):} Administers tax collection
        and enforces tax law. Revenue collection is the CRA's core function;
        it does not make tax policy, which remains with Finance Canada.
  \item \textbf{Employment and Social Development Canada (ESDC):} Administers
        transfer programs including Employment Insurance (EI), the Canada
        Pension Plan (CPP), the Canada Child Benefit (CCB), and Old Age
        Security (OAS/GIS). These programs collectively transfer hundreds
        of billions of dollars per year from taxpayers to households.
\end{itemize}

%----------------------------------------------------------
\subsection{Statistics Canada}
%----------------------------------------------------------

\term{Statistics Canada} (\textit{Statistique Canada}) is the national
statistical agency, established under the \textit{Statistics Act}. It
is the primary data source for the analysis conducted throughout this
textbook. Its publications include:

\begin{itemize}
  \item \textbf{Consumer Price Index (CPI):} Published monthly; the primary
        measure of inflation (Table 18-10-0004-01). Analysed in depth in
        Chapter~\ref{ch:inflation}.
  \item \textbf{Labour Force Survey (LFS):} Published monthly; measures the
        employment rate, unemployment rate, and labour force participation
        rate (Table 14-10-0287-01). Central to Chapter~\ref{ch:labour}.
  \item \textbf{National Income and Expenditure Accounts:} Quarterly GDP by
        expenditure, income, and industry (Table 36-10-0104-01). The
        foundation for Chapters~\ref{ch:introduction} through \ref{ch:fiscal}.
  \item \textbf{Survey of Consumer Finances / Canadian Income Survey:}
        Annual data on household income distribution, poverty, and inequality
        (Table 11-10-0190-01). Central to Chapter~\ref{ch:inequality}.
  \item \textbf{Census of Population:} Every five years; provides
        comprehensive demographic, labour market, and income data. The 2021
        Census is heavily used in Chapters~\ref{ch:inequality}
        and \ref{ch:immigration}.
\end{itemize}

Statistics Canada operates with formal independence from political direction
in its data collection and methodological decisions. The integrity of its
data is a fundamental public good: without reliable statistics, evidence-based
policymaking is impossible.

%----------------------------------------------------------
\subsection{Canada Mortgage and Housing Corporation (CMHC)}
%----------------------------------------------------------

The \term{Canada Mortgage and Housing Corporation} (CMHC) is a federal
Crown corporation with two distinct functions. First, it operates a
\term{mortgage insurance} program: Canadians who purchase a home with
a down payment of less than 20\% must purchase CMHC mortgage insurance,
which protects the lender against default. This insurance program
allows lenders to extend mortgages to buyers with limited savings, making
homeownership accessible to a broader population. Second, CMHC is the
primary source of Canadian housing market data --- including housing starts,
rental vacancy rates, affordability indices, and housing market outlooks ---
and it is the implementing agency for the federal government's housing programs.
CMHC is central to the analysis in Chapter~\ref{ch:housing}.

%----------------------------------------------------------
\subsection{Provincial Governments}
%----------------------------------------------------------

In Canada's federal system, provincial governments bear primary responsibility
for several of the largest areas of economic policy:

\begin{itemize}
  \item \textbf{Healthcare:} Provinces operate and fund the publicly administered
        insurance plans mandated by the \textit{Canada Health Act}. Healthcare
        is typically the single largest provincial spending item --- consuming
        35--45\% of provincial budgets.
  \item \textbf{Education:} Provinces regulate and largely fund education from
        kindergarten through post-secondary.
  \item \textbf{Labour market regulation:} Minimum wages, occupational safety,
        labour relations, and most employment standards are set provincially.
        This is why minimum wages vary from \$13.85 in some Atlantic provinces
        to \$17.40 in British Columbia.
  \item \textbf{Natural resources:} Under the Constitution Act, provinces own
        the natural resources within their borders. Provincial royalty regimes
        for oil, gas, and minerals are major determinants of resource sector
        economics.
  \item \textbf{Housing regulation:} Zoning, land use planning, and building
        codes are primarily municipal functions delegated by provinces. This
        makes provinces central players in the housing affordability crisis
        analysed in Chapter~\ref{ch:housing}.
\end{itemize}

\begin{table}[H]
\centering
\caption{Canada's Major Economic Institutions: Mandate, Tools, and Independence}
\label{tab:institutions}
\renewcommand{\arraystretch}{1.35}
\begin{tabular}{p{3.1cm}p{3.6cm}p{2.8cm}p{2.8cm}}
\toprule
\textbf{Institution} & \textbf{Primary Mandate} & \textbf{Main Tools}
                     & \textbf{Independence Level} \\
\midrule
Bank of Canada
  & Low and stable inflation (2\% target); financial stability
  & Overnight rate; quantitative easing; financial regulation
  & High --- Crown corporation, arm's length from government \\
\addlinespace
Department of Finance Canada
  & Fiscal policy; budget; taxation
  & Government spending; tax rates; transfer programs
  & Low --- elected minister, accountable to Parliament \\
\addlinespace
Statistics Canada
  & Collect, produce, and disseminate national statistics
  & Surveys; censuses; administrative data
  & High --- statistical independence protected by law \\
\addlinespace
CMHC
  & Housing finance; mortgage insurance; housing policy
  & Mortgage insurance; housing programs; market data
  & Moderate --- Crown corporation with federal mandate \\
\addlinespace
Parliamentary Budget Officer
  & Independent fiscal and economic analysis
  & Economic modelling; budget assessments; reports
  & High --- independent officer of Parliament \\
\addlinespace
Provincial Governments
  & Healthcare; education; labour; natural resources
  & Spending; regulation; taxation; royalties
  & Full --- constitutionally sovereign within their jurisdiction \\
\bottomrule
\end{tabular}
\source{Bank of Canada Act; Statistics Act; Canada Mortgage and Housing
Corporation Act; Constitution Act, 1867.}
\end{table}

%==========================================================
\section{The Policy Toolkit: Monetary and Fiscal Policy}
\label{sec:policy_toolkit}
%==========================================================

Two broad categories of macroeconomic policy tools are the primary instruments
through which governments and central banks attempt to influence the
performance of the economy. Understanding the distinction between them ---
and the interaction between them --- is foundational for every topic in
this course.

%----------------------------------------------------------
\subsection{Monetary Policy}
%----------------------------------------------------------

\begin{defbox}{Monetary Policy}
\textbf{Monetary policy} refers to the actions taken by a central bank
to influence the money supply, interest rates, and credit conditions,
with the ultimate goals of maintaining price stability and supporting
sustainable economic growth.

\textit{Economic intuition:} When the central bank raises its policy
interest rate, borrowing becomes more expensive throughout the economy.
Households take out smaller mortgages. Firms postpone capital investment.
The exchange rate tends to appreciate, making imports cheaper and exports
less competitive. These effects reduce aggregate demand and, over time,
inflation. The reverse occurs when the Bank cuts rates.

\textit{Key constraint --- the zero lower bound:} The policy rate cannot
be reduced significantly below zero (depositors would withdraw cash and
hold it physically rather than accept negative returns on deposits). This
creates an asymmetry: the Bank can always raise rates to fight inflation,
but its ability to cut rates to fight recession is constrained at zero.
After 2008 and again in 2020, the Bank cut its rate to 0.25\% --- its
effective lower bound --- and was forced to rely on other tools (forward
guidance and, to a limited extent, quantitative easing) to provide
additional stimulus.

\textit{Canadian example:} The Bank cut the overnight rate from 1.75\%
to 0.25\% in March 2020 in response to COVID-19, then raised it to 5.00\%
between March 2022 and July 2023 to fight inflation --- the full range of
monetary policy in action within a four-year period.
\end{defbox}

The primary transmission channels through which a Bank of Canada interest
rate change affects the broader economy are:

\begin{enumerate}
  \item \textbf{Credit cost channel:} Higher overnight rates are passed
        through to mortgage rates, car loan rates, and business borrowing
        costs, reducing household and business spending.
  \item \textbf{Asset price channel:} Higher interest rates reduce the
        present value of future asset cash flows, lowering stock prices
        and real estate values, which reduces household wealth and
        consumption (the ``wealth effect'').
  \item \textbf{Exchange rate channel:} Higher Canadian interest rates
        attract foreign capital seeking better returns, increasing demand
        for Canadian dollars and appreciating the exchange rate. A stronger
        Canadian dollar makes imports cheaper (reducing inflation) and
        exports more expensive (reducing demand for Canadian goods abroad).
  \item \textbf{Expectations channel:} Perhaps most importantly, when
        the Bank credibly commits to fighting inflation, households and
        firms adjust their price-setting and wage demands downward,
        reducing inflationary momentum even before interest rates fully
        transmit through the system.
\end{enumerate}

%----------------------------------------------------------
\subsection{Fiscal Policy}
%----------------------------------------------------------

\begin{defbox}{Fiscal Policy}
\textbf{Fiscal policy} refers to the use of government spending and
taxation to influence macroeconomic conditions.

\textit{Economic intuition:} When aggregate demand is insufficient to
support full employment, the government can increase its own spending
(shifting the aggregate demand curve rightward) or cut taxes (increasing
household disposable income and thus consumption). Conversely, when
the economy is overheating, fiscal restraint --- cutting spending or
raising taxes --- can reduce inflationary pressure.

\textit{Key constraint --- the deficit and debt:} When the government
spends more than it collects in taxes, it runs a \term{budget deficit},
which must be financed by borrowing. The accumulation of past deficits
is the \term{national debt}. Sustained high deficits raise the stock
of debt relative to GDP, increasing interest payments and potentially
crowding out private investment. The fiscal arithmetic of debt
sustainability is explored fully in Chapter~\ref{ch:fiscal}.

\textit{Canadian example:} In 2020--21, the federal government ran a
deficit of \$327.7~billion --- 14.9\% of GDP --- as it deployed the
Canada Emergency Response Benefit (CERB), the Canada Emergency Wage
Subsidy (CEWS), and other programs to protect household incomes during
the pandemic shutdown. This was the largest single-year deficit in
Canadian history. It was deliberate and effective in preventing a
financial collapse, but it also raised Canada's federal debt-to-GDP
ratio from approximately 31\% to 49\% within two years.
\end{defbox}

%----------------------------------------------------------
\subsection{The Interaction Between Monetary and Fiscal Policy}
%----------------------------------------------------------

Monetary and fiscal policy are operated by different institutions with
different mandates, but they interact in important ways. When fiscal
policy is highly expansionary --- as it was in 2020--21 --- it can
increase aggregate demand and put upward pressure on inflation, requiring
a tighter monetary policy response. Conversely, when monetary policy is
very tight --- as it was in 2022--23 --- higher interest rates increase
the government's debt servicing costs, creating fiscal pressure.

The period 2020--2023 illustrates this interaction vividly. Emergency
fiscal spending in 2020--21 was appropriate and necessary. But the
combination of massive fiscal stimulus and near-zero interest rates
through 2021 contributed to the demand-side component of the inflation
surge that the Bank of Canada then had to combat aggressively in 2022--23.
Understanding this interaction is essential for evaluating both the
monetary policy response (Chapter~\ref{ch:inflation}) and the fiscal
policy choices (Chapter~\ref{ch:fiscal}).

\begin{figure}[H]
\centering
\begin{tikzpicture}[
  node distance=1.6cm and 2.8cm,
  mainbox/.style={rectangle, rounded corners=5pt, draw=unbcgreen,
                  fill=unbcgreenlight, minimum width=3.2cm,
                  minimum height=1.1cm, align=center, font=\small\bfseries},
  midbox/.style={rectangle, rounded corners=4pt, draw=policyblue,
                 fill=policyblulight, minimum width=3.0cm,
                 minimum height=0.95cm, align=center, font=\small},
  finalbox/.style={rectangle, rounded corners=5pt, draw=unbcgreen,
                   fill=unbcgreenmid, minimum width=7cm,
                   minimum height=1.0cm, align=center, font=\small\bfseries},
  arrow/.style={-Stealth, thick, color=unbcgreen!80}
]
\node[mainbox] (boc)     {Bank of Canada\\(Monetary Policy)};
\node[mainbox, right=3.6cm of boc] (fin) {Dept.\ of Finance\\(Fiscal Policy)};

\node[midbox, below=1.6cm of boc] (ir)   {Overnight Rate\\\& Credit Costs};
\node[midbox, below=1.6cm of fin] (gt)   {Spending, Taxes\\\& Transfers};

\node[midbox, below=1.6cm of ir]  (inv)  {Investment,\\ Consumption,\\ Exchange Rate};
\node[midbox, below=1.6cm of gt]  (ad2)  {Household\\ Disposable\\ Income};

\node[finalbox, below=4.8cm of $(boc)!0.5!(fin)$] (out)
     {Aggregate Demand $\longrightarrow$ GDP Growth, Employment, Inflation};

\draw[arrow] (boc) -- (ir);
\draw[arrow] (fin) -- (gt);
\draw[arrow] (ir)  -- (inv);
\draw[arrow] (gt)  -- (ad2);
\draw[arrow] (inv) -- (out);
\draw[arrow] (ad2) -- (out);

\draw[dashed, color=gray, thick, <->]
  ([xshift=0.5cm]boc.east) -- ([xshift=-0.5cm]fin.west)
  node[midway, above, font=\footnotesize, color=gray]
  {Coordination\,/\,Conflict};

\node[above=0.4cm of boc, font=\small\bfseries\color{unbcgreen}]
  {CANADA'S MACROECONOMIC POLICY ARCHITECTURE};
\end{tikzpicture}
\caption{Canada's Macroeconomic Policy Architecture.
Monetary policy (left branch) and fiscal policy (right branch) each
affect aggregate demand through distinct channels, ultimately influencing
GDP growth, employment, and inflation. The dashed bidirectional arrow
indicates that the two policy levers can work in the same direction
(as during the COVID stimulus) or in opposing directions (as in 2022--23,
when fiscal policy remained expansionary even as the Bank of Canada
was tightening). Understanding this interaction is one of the central
analytical tasks of contemporary macroeconomics.}
\label{fig:policy_framework}
\end{figure}

\begin{policydebate}{Should the Bank of Canada Be Independent from the Federal Government?}

\textbf{The Issue:} Central bank independence --- the insulation of
monetary policy decisions from direct political control --- is a cornerstone
of modern macroeconomic governance. But independence is not constitutionally
guaranteed in Canada. The \textit{Bank of Canada Act} allows the federal
government, in exceptional circumstances, to direct the Bank. Should this
arm's-length relationship be strengthened, weakened, or left as is?

\textbf{The Case for Independence:}
Economic theory and empirical evidence both support central bank independence.
The fundamental problem is \term{time inconsistency}: elected governments
have incentives to expand the economy before elections, even at the cost of
higher future inflation. If markets expect governments to exploit monetary
policy for electoral gain, they will build that expectation into wages and
prices, requiring permanently higher unemployment to achieve any given
inflation rate. An independent central bank can credibly commit to price
stability in a way that elected governments cannot --- and that credibility
is itself enormously valuable. Cross-country evidence consistently shows
that countries with more independent central banks have lower average
inflation without worse unemployment or growth outcomes
\citep[see][]{alesina1993,cukierman1992}.

\textbf{The Case Against Full Independence:}
Critics argue that central bank independence raises a \term{democratic
accountability} problem. When the Bank of Canada raises interest rates
by 475 basis points in 16 months --- dramatically increasing mortgage
costs for millions of Canadians --- this is a consequential distributional
decision that arguably deserves more democratic oversight than it currently
receives. Furthermore, the separation of monetary and fiscal policy can
make macroeconomic management incoherent: if the elected government wishes
to support employment and the central bank is tightening to control
inflation, the two policies work against each other, potentially prolonging
both the high interest rates and the high unemployment.

\textbf{The Canadian Position:}
Canada's current arrangement is a hybrid. The Bank of Canada is independent
in its day-to-day operations, but its inflation target is set jointly
with the Government of Canada in a periodic agreement (most recently renewed
in 2021). This ``joint mandate'' approach attempts to combine the credibility
benefits of independence with some degree of democratic input into the
overarching framework.

\textbf{Your Assessment:}
After completing Chapter~\ref{ch:inflation}, you will be equipped to evaluate
the Bank's rate-hike cycle on its merits. Did the Bank move too aggressively?
Too slowly? At the right pace? These are exactly the kinds of questions where
the distinction between positive analysis (what happened and what were the
effects?) and normative judgement (was this the right policy?) matters most.
\end{policydebate}

%==========================================================
\section{How Economists Think}
\label{sec:how_economists_think}
%==========================================================

Developing the analytical habits of economic reasoning takes time and
deliberate practice. This section introduces the core tools that you will
use throughout the course.

%----------------------------------------------------------
\subsection{Economic Models and Assumptions}
%----------------------------------------------------------

Economists use \term{models} --- carefully constructed simplifications of
reality --- to generate predictions and organize thinking about complex
phenomena. A model deliberately omits details in order to focus on the
mechanisms that most matter for a particular question.

The \term{supply and demand model}, for example, abstracts away from the
unique characteristics of millions of individual buyers and sellers and
focuses instead on the aggregate relationship between price and quantity.
This abstraction is a feature, not a limitation: it allows economists to
make clear, testable predictions about how markets respond to changes in
supply, demand, and policy.

Every model rests on assumptions. Evaluating a model requires understanding
which assumptions are load-bearing (if they are wrong, the model's
predictions fail) and which are harmless simplifications. The assumption
that economic agents ``respond to incentives'' is load-bearing for almost
all of economics. The assumption of perfect information is often harmless
as a first approximation but becomes crucial in analyses of insurance,
financial markets, and healthcare.

\begin{thinkbox}
A map is a model. A road map of Canada omits trees, buildings, and
terrain because those details are irrelevant to the task of navigating
from Prince George to Edmonton. A topographic map of the same region
omits roads but captures elevation, because that information matters
for a hiker. Neither map is ``more correct'' than the other --- they are
optimized for different questions. Economic models work the same way.
A model that is useful for understanding inflation may be useless for
understanding income inequality. The skill is in choosing the right
model for the question at hand.
\end{thinkbox}

%----------------------------------------------------------
\subsection{Positive and Normative Economics}
%----------------------------------------------------------

One of the most important distinctions in economic methodology is between
\term{positive economics} and \term{normative economics}.

\begin{defbox}{Positive and Normative Economics}
\textbf{Positive economics} describes and explains economic phenomena
as they are, without making value judgements. It answers ``what is''
or ``what will happen if.'' Positive claims are, in principle, testable
against empirical evidence.

\textbf{Normative economics} makes value judgements about what
\textit{should} happen. It answers ``what ought to be.'' Normative
claims cannot be resolved by evidence alone --- they also depend on
values and priorities.

\textbf{Examples in the Canadian context:}
\begin{itemize}
  \item \textit{Positive:} ``Canada's CPI inflation peaked at 8.1\% in
        June 2022.'' [Observable fact.]
  \item \textit{Positive:} ``A 25\% US tariff on Canadian auto exports
        would reduce Canadian GDP by an estimated 0.3--0.7 percentage
        points.'' [Testable empirical claim, though uncertain.]
  \item \textit{Normative:} ``The federal carbon tax should be increased
        to \$250 per tonne.'' [Value judgement about the right balance
        between climate and economic costs.]
  \item \textit{Normative:} ``Canada should prioritize reducing inequality
        over maximizing GDP growth.'' [Value judgement about priorities.]
\end{itemize}

Note: many policy debates conflate positive and normative claims. Part
of the discipline of economics is keeping them separate.
\end{defbox}

Policy debates are often confusing precisely because positive and normative
claims are mixed together. When a politician argues that the carbon tax
``kills jobs'' (positive claim --- testable) and ``should be abolished''
(normative claim --- based on values), they are making two distinct types
of argument that require two distinct types of response. Throughout this
textbook, we will be explicit about which claims are positive (and what
the evidence says) and which are normative (and what the competing values are).

%----------------------------------------------------------
\subsection{Opportunity Cost and Trade-offs}
%----------------------------------------------------------

A fundamental principle of economic reasoning is that \textit{all choices
involve trade-offs}. Resources --- time, money, productive capacity, policy
attention --- allocated to one purpose cannot simultaneously be used for
another. The \term{opportunity cost} of any choice is the value of the
next-best alternative foregone.

\begin{equation}
  \text{Opportunity Cost of Choice A} = \text{Value of the best foregone alternative to A}
  \label{eq:opportunity_cost}
\end{equation}

Opportunity cost is not the same as monetary cost. When a university student
works part-time during the school year, the opportunity cost of those hours
worked is not just the forgone leisure but the forgone study time and the
potential impact on grades. When the federal government allocates
\$15~billion to the National Housing Strategy, the opportunity cost is not
just the \$15~billion in budget space but whatever the next-best use of
that money would have been --- whether tax cuts, healthcare investment,
debt repayment, or climate adaptation spending.

\begin{canexample}{The Opportunity Cost of Canada's COVID Fiscal Response}
In 2020--21, the federal government spent approximately \$350 billion in
emergency COVID-19 support. This was financed by borrowing, raising the
federal debt-to-GDP ratio by nearly 18 percentage points. The opportunity
cost of this spending was not only the future interest payments
(approximately \$43 billion per year by 2024--25), but also the fiscal
space foregone for other priorities: a national pharmacare program, a
substantial climate investment fund, or a faster post-pandemic deficit
reduction path. The emergency spending was arguably necessary and worth
its cost --- but a complete analysis must acknowledge the trade-offs.
Chapter~\ref{ch:fiscal} revisits this question with full fiscal data.
\end{canexample}

%----------------------------------------------------------
\subsection{Marginal Analysis}
%----------------------------------------------------------

Economists typically analyse decisions not as binary choices (do this or do
not do this) but at the \term{margin}: how much of something should be done?
\term{Marginal cost} (MC) is the additional cost of one more unit of
an activity. \term{Marginal benefit} (MB) is the additional benefit.
The economically optimal decision equates the two:

\begin{equation}
  \text{Optimal decision:} \quad MB = MC
  \label{eq:mb_mc}
\end{equation}

The marginal analysis framework applies broadly. Should the Bank of Canada
raise interest rates another 25 basis points? The marginal benefit is a
further reduction in inflation; the marginal cost is a further increase
in mortgage payments and a risk of over-tightening into recession. Should
Canada raise the carbon tax from \$65 to \$80 per tonne? The marginal
benefit is additional emissions reduction; the marginal cost is additional
burden on households and firms. In both cases, the right question is not
``should we do more?'' in the abstract, but ``does the marginal benefit
exceed the marginal cost at this particular point?''

%----------------------------------------------------------
\subsection{Correlation and Causation}
%----------------------------------------------------------

Perhaps the most important methodological lesson in empirical economics is
the distinction between \term{correlation} and \term{causation}. Two
variables are correlated if they tend to move together. But correlation
does not establish causation: variable A may cause B, B may cause A, or
a third variable C may cause both A and B.

\begin{thinkbox}
Ice cream sales and drowning deaths are positively correlated --- both rise
in summer. Does eating ice cream cause drowning? Of course not: the third
variable, warm weather, drives both. Identifying this kind of spurious
correlation is easy. In economics, the confounding variables are often
far less obvious.

Consider: provinces with higher minimum wages tend to have higher
unemployment rates. Does this prove that minimum wages cause unemployment?
Not necessarily --- provinces with historically weaker labour markets may
have adopted higher minimum wages as a compensatory policy, or other
factors (industry mix, demographics) may explain both. Establishing
\textit{causal} effects from observational data requires careful
research design, and the results are often contested. We will encounter
this challenge repeatedly throughout the course.
\end{thinkbox}

%----------------------------------------------------------
\subsection{Efficiency and Equity}
%----------------------------------------------------------

A final pair of concepts that runs through every chapter is the distinction
between \term{economic efficiency} and \term{economic equity} (or fairness).
An allocation of resources is \term{efficient} if it is impossible to make
someone better off without making someone else worse off --- a
\term{Pareto optimal} state. Markets, under the right conditions, tend to
produce efficient outcomes. But efficient outcomes need not be equitable:
a distribution in which one person owns everything and everyone else
owns nothing is technically Pareto optimal (you cannot help the poor
without taking from the rich).

Many of the most important policy debates in this course --- should Canada
have a carbon tax? should it have rent control? should it have a wealth tax?
--- involve precisely this tension. A policy that increases efficiency
(total GDP) may worsen equity (distribution of that GDP), or vice versa.
Evaluating these trade-offs requires clear thinking about both what the
evidence says (positive) and what values should guide our priorities
(normative).

%==========================================================
\section{Canada in International Context}
\label{sec:international_context}
%==========================================================

Understanding Canada's contemporary economic situation requires situating
it within the broader international environment. Canada's economic performance
is shaped by global forces --- commodity prices, US monetary policy, global
trade flows, and climate commitments --- that no Canadian institution can
fully control.

\begin{figure}[H]
\centering
\begin{tikzpicture}
\begin{axis}[
  width=14cm, height=7.5cm,
  xlabel={Year},
  ylabel={Real GDP Growth Rate (\%)},
  xlabel style={font=\small},
  ylabel style={font=\small},
  xmin=1999, xmax=2027,
  ymin=-7, ymax=7,
  xtick={2000,2002,2004,2006,2008,2010,2012,2014,2016,2018,2020,2022,2024,2026},
  xticklabel style={rotate=40, anchor=east, font=\footnotesize},
  ytick={-6,-4,-2,0,2,4,6},
  yticklabel style={font=\small},
  ymajorgrids=true,
  grid style={dashed, gray!25},
  axis line style={color=gray!60},
  tick style={color=gray!60},
  legend pos=north east,
  legend style={font=\small},
  clip=false
]
% Reference line at 0%
\draw[dashed, gray, thick] (axis cs:1999,0) -- (axis cs:2027,0);
% Shaded recession bands
\fill[gray!15] (axis cs:2008.5,-7) rectangle (axis cs:2009.5,7);
\fill[gray!15] (axis cs:2019.8,-7) rectangle (axis cs:2020.8,7);
% Real GDP growth line
\addplot[unbcgreen, thick, mark=*, mark size=1.5pt] coordinates {
  (2000, 5.2)(2001, 1.8)(2002, 3.0)(2003, 1.9)(2004, 3.1)
  (2005, 3.2)(2006, 2.6)(2007, 2.1)(2008, 1.0)(2009,-2.9)
  (2010, 3.1)(2011, 3.1)(2012, 1.7)(2013, 2.5)(2014, 2.9)
  (2015, 0.7)(2016, 1.0)(2017, 3.0)(2018, 2.4)(2019, 1.9)
  (2020,-5.2)(2021, 5.0)(2022, 3.4)(2023, 1.2)(2024, 1.5)
  (2025, 1.3)(2026, 1.8)
};
% Potential growth band
\addplot[dataorange, dashed, thick] coordinates {
  (2000,2.0)(2026,2.0)
} node[right, font=\footnotesize, xshift=2pt] {Potential growth ($\approx$2\%)};

% Annotations
\node[font=\footnotesize, align=center, text=gray] at (axis cs:2009,-4.5)
  {Financial\\Crisis\\$-$2.9\%};
\node[font=\footnotesize, align=center, text=gray] at (axis cs:2015.5,-1.5)
  {Oil price\\shock};
\node[font=\footnotesize, align=center, text=gray] at (axis cs:2020,-6.3)
  {COVID-19\\$-$5.2\%};
\node[font=\footnotesize, align=center, text=unbcgreen] at (axis cs:2021,6.2)
  {Rebound\\+5.0\%};
\addlegendentry{Canada Real GDP Growth}
\end{axis}
\end{tikzpicture}
\caption{Canada's Real GDP Growth Rate, 2000--2026. The shaded bands mark
recessionary years. Three episodes stand out: the 2009 financial crisis
contraction ($-$2.9\%), driven by the global banking collapse and the
near-simultaneous US recession; the 2015--16 slowdown driven by the oil
price collapse to below \$30/barrel (US); and the 2020 COVID-19 collapse
($-$5.2\%), the steepest in peacetime Canadian history. The dashed orange
line at 2\% represents Canada's approximate \textit{potential} growth rate ---
the rate consistent with stable unemployment and inflation. Years above this
line indicate excess demand; years below it indicate slack.}
\label{fig:gdp_growth}
\source{Statistics Canada, Table 36-10-0104-01, Chained (2017) dollars.}
\end{figure}

\begin{table}[H]
\centering
\caption{Key Canadian Macroeconomic Indicators, 2015--2026}
\label{tab:macro_indicators}
\renewcommand{\arraystretch}{1.30}
\begin{tabular}{lrrrrrr}
\toprule
\textbf{Year} & \textbf{Real GDP} & \textbf{CPI} & \textbf{Unemp.}
              & \textbf{Bank} & \textbf{Fed.\ Balance} & \textbf{Fed.\ Debt} \\
              & \textbf{Growth (\%)} & \textbf{Infl.\ (\%)} & \textbf{Rate (\%)}
              & \textbf{Rate (\%)} & \textbf{(\$B)} & \textbf{(\% GDP)} \\
\midrule
2015 & 0.7  & 1.1  & 6.9 & 0.50 & $-$2.9   & 31.0 \\
2016 & 1.0  & 1.4  & 7.0 & 0.50 & $-$17.8  & 31.5 \\
2017 & 3.0  & 1.6  & 6.3 & 1.00 & $-$19.0  & 30.8 \\
2018 & 2.4  & 2.3  & 5.8 & 1.75 & $-$14.0  & 30.1 \\
2019 & 1.9  & 1.9  & 5.7 & 1.75 & $-$39.4  & 30.8 \\
2020 & $-$5.2 & 0.7 & 9.5 & 0.25 & $-$327.7 & 49.1 \\
2021 & 5.0  & 3.4  & 7.5 & 0.25 & $-$113.8 & 45.2 \\
2022 & 3.4  & 6.8  & 5.3 & 4.25 & $-$36.0  & 42.4 \\
2023 & 1.2  & 3.9  & 5.4 & 5.00 & $-$40.0  & 42.0 \\
2024 & 1.5  & 2.6  & 6.3 & 4.25 & $-$46.8  & 42.3 \\
2025 & 1.3  & 2.1  & 6.7 & 3.00 & $-$48.0  & 42.5 \\
2026 & 1.8  & 2.0  & 6.5 & 2.75 & $-$30.0  & 42.0 \\
\bottomrule
\end{tabular}
\source{Statistics Canada, Tables 36-10-0104-01, 18-10-0004-01,
14-10-0287-01; Bank of Canada Key Interest Rates; Department of Finance
Canada, Fiscal Reference Tables, 2026. Values for 2025--2026 are
preliminary or projected. Federal balance figures in nominal Canadian
dollars; negative values indicate a deficit.}
\end{table}

\noindent Table~\ref{tab:macro_indicators} is the single most important
reference table in this textbook. It captures the macroeconomic story of
Canada across 12 years in six dimensions. Study it carefully. Notice that
2020 is the anomaly year on every dimension simultaneously: growth collapses,
inflation vanishes, unemployment spikes, the Bank rate hits its floor, and
the budget deficit reaches a historic extreme. Notice that 2022 is the
year of the great inflation --- 6.8\%, the highest since 1991. Notice that
the Bank rate and CPI inflation move together through 2022--2024, reflecting
the monetary policy response. And notice that the federal debt-to-GDP ratio
rose by nearly 18 percentage points in a single year (2019 to 2020) and
has not fully recovered --- a fiscal legacy that Chapter~\ref{ch:fiscal}
examines in depth.

%==========================================================
\section{Case Studies}
\label{sec:ch1_casestudies}
%==========================================================

%----------------------------------------------------------
\begin{casestudy}{COVID-19's Economic Shock and Canada's Recovery}

\textbf{Background:} On March 11, 2020, the World Health Organization
declared COVID-19 a global pandemic. Within days, Canadian provinces
ordered the closure of non-essential businesses, schools, and public
spaces. Economic activity collapsed at an unprecedented speed.

\textbf{The Economic Shock:} In April 2020, Canada lost approximately
2~million jobs --- more than were lost in the entire 2008--09 recession.
GDP fell at an annualized rate exceeding 40\% in the second quarter of 2020.
The unemployment rate reached 13.7\% in May 2020. These numbers, while
startling, understate the disruption: millions of additional Canadians
were employed but working zero hours (classified as employed but absent
by the LFS methodology), and the labour force participation rate fell
sharply as many stopped looking for work.

\textbf{Policy Response:} The federal government deployed the Canada Emergency
Response Benefit (CERB), which paid \$2,000 per month to approximately
9 million workers who lost employment income --- more than one in four
Canadian workers. The Canada Emergency Wage Subsidy (CEWS) paid up to
75\% of the wages of employees at firms whose revenues had fallen, at a
total cost exceeding \$100 billion. The Bank of Canada cut its policy
rate to 0.25\% and committed to holding it there until the recovery was
``well underway.'' Combined, the fiscal and monetary response was the
most aggressive in Canadian peacetime history.

\textbf{Recovery:} Canada's recovery was faster than most projections.
By early 2022, employment had fully recovered to pre-pandemic levels.
GDP recovered its lost ground by early 2022 as well. But the recovery
was uneven: high-contact service sectors (restaurants, hotels, live
entertainment) recovered more slowly; women, racialized Canadians, and
low-wage workers experienced longer-lasting earnings impacts; and the
combination of massive fiscal stimulus, low interest rates, and supply
chain disruptions created the inflationary surge of 2022 that reshaped
the economic environment for years afterward.

\textbf{Evaluation:} The COVID episode illustrates both the power and the
limits of macroeconomic policy. The fiscal and monetary response prevented
a financial collapse and protected millions of households. But it also
contributed to inflationary dynamics that the Bank of Canada then had to
combat at significant cost to mortgage-holding households. Understanding
this sequence of events --- pandemic shock, policy response, inflationary
consequence, monetary tightening, housing market correction --- provides
the narrative architecture for Chapters~\ref{ch:inflation} through
\ref{ch:housing}.

\textit{Data reference:} Statistics Canada, Table 14-10-0287-01 (employment);
Table 36-10-0104-01 (GDP); Parliamentary Budget Officer, COVID-19
Fiscal Analysis, 2021.
\end{casestudy}

%----------------------------------------------------------
\begin{casestudy}{Northern British Columbia as a Regional Economy}

\textbf{Background:} The University of Northern British Columbia sits in
Prince George --- the largest city in northern BC, with a population of
approximately 75,000. The regional economy of northern BC illustrates,
in compressed form, many of the structural features and contemporary
challenges of the Canadian economy as a whole.

\textbf{Economic Structure:} Northern BC's economy is built on natural
resources: forestry, mining (gold, copper, coal), and increasingly, natural
gas from the Montney formation and the Peace River region. The completion
of the Trans Mountain Pipeline Expansion (2024) and the LNG Canada facility
near Kitimat (under construction) represent billion-dollar investments that
will shift the regional resource mix toward hydrocarbons. The forest sector,
historically the dominant employer, has been reshaped by mill closures,
automation, and the mountain pine beetle epidemic that devastated the lodgepole
pine forests of the 1990s and 2000s.

\textbf{Labour Market Features:} Northern BC labour markets are characterized
by: above-average wages in resource sectors (unionized sawmill and mine
workers earn substantially more than provincial minimum wage); persistent
labour shortages in healthcare and social services (geographic remoteness
makes recruitment difficult); a significant Indigenous population (First
Nations account for approximately 16\% of the regional population) whose
economic participation and outcomes reflect the historical legacy of
dispossession analysed in Chapter~\ref{ch:inequality}; and a mismatch
between the skills demanded by resource employers and those produced by
the local education system.

\textbf{Contemporary Challenges:} The three most acute economic challenges
in northern BC today mirror those of the national economy: a housing market
that has become significantly less affordable as in-migration from
Metro Vancouver and Alberta has driven up prices; a healthcare system
under fiscal pressure (Northern Health Authority); and the question of
how a resource-dependent regional economy navigates the clean energy
transition while maintaining employment and fiscal revenues. These regional
challenges appear in the course not as peripheral ``local colour'' but as
concrete examples through which national economic forces become tangible.
\end{casestudy}

%----------------------------------------------------------
\begin{casestudy}{Canada's 2022 Inflation Surge: A Turning Point}

\textbf{Background:} Canada had not experienced inflation above 5\% since
1991. A generation of economists, policymakers, and households had come
of age in a low-inflation environment. The 2021--2022 surge was therefore
not merely an economic event but a psychological and institutional stress test.

\textbf{What Happened:} Consumer price inflation rose from 0.7\% in 2020
(deflated by pandemic-related demand collapse) to 3.4\% in 2021 and then
to a peak of 8.1\% in June 2022. By that point, food prices were rising
at over 10\% year-over-year. Shelter costs were rising at over 7\%.
Energy prices --- driven by Russian sanctions following the Ukraine invasion ---
were rising at over 30\% for gasoline and over 20\% for natural gas.

\textbf{Why It Happened:} The inflation surge had multiple, simultaneous causes.
On the demand side: the COVID fiscal stimulus had placed over \$350 billion
into Canadian household bank accounts and business balance sheets, creating
excess demand when supply chains could not expand fast enough to meet it.
Low interest rates through 2021 added further fuel. On the supply side:
global shipping costs rose tenfold; semiconductor shortages constrained
automobile and electronics production; pandemic-related labour shortages
reduced output in food processing, transportation, and construction. The
Russian invasion of Ukraine (February 2022) added an energy and food
commodity price shock on top of already-elevated inflation.

\textbf{Policy Response:} The Bank of Canada raised its policy rate 10 times
between March 2022 and July 2023, from 0.25\% to 5.00\%. By end-2023,
inflation had fallen to 3.4\%; by end-2024, it was approaching 2\%.
But the rate hikes came at significant cost: mortgage payments for variable-rate
holders increased by hundreds of dollars per month; housing affordability
deteriorated further; and economic growth slowed sharply.

\textbf{Lesson:} The 2022 episode illustrates that macroeconomic stability
cannot be taken for granted. Institutional frameworks (like the inflation-targeting
regime) that functioned smoothly for three decades can be tested by the right
combination of shocks. It also illustrates the limits of monetary policy:
the Bank of Canada can raise rates to reduce demand, but it cannot directly
fix supply chain disruptions, energy prices determined in global markets,
or geopolitical shocks. These limits are a recurring theme in
Chapter~\ref{ch:inflation}.
\end{casestudy}

%==========================================================
\section{Economics in Everyday Life}
%==========================================================

\begin{everydaylife}{From the Gas Station to the Central Bank: Tracing Economic Forces in Your Daily Life}
The next time you fill up your car, consider the chain of economic
forces determining the number on the pump.

The \textbf{global oil price} (set in markets in Houston and Amsterdam)
determines the base cost of crude. Canada's \textbf{exchange rate} ---
the value of the Canadian dollar relative to the US dollar --- translates
that global price into Canadian dollars: a 10-cent drop in the loonie
raises the price of a barrel of oil (priced in USD) by approximately
7--8\% in Canadian terms. The \textbf{Bank of Canada's interest rate
decisions} affect the exchange rate by influencing capital flows into and
out of Canada. Federal and provincial \textbf{excise taxes} add
approximately 40--50 cents per litre at the pump --- of which about
17 cents is the federal carbon levy (at the 2024 rate of \$65 per tonne).
\textbf{Refinery margins}, influenced by the number and capacity of
Canadian refineries (a market structure question), add another layer of cost.

Every chapter in this textbook touches on one of these forces. Chapter~\ref{ch:inflation}
examines how energy prices drive CPI inflation. Chapter~\ref{ch:trade}
examines the exchange rate. Chapter~\ref{ch:fiscal} examines excise taxes.
Chapter~\ref{ch:climate} examines the carbon levy. By the end of this course,
the number on that gas pump will be transparent to you in a way it is not
for most Canadians --- not because you will agree with every policy that
produced it, but because you will understand the economic forces behind each
component.
\end{everydaylife}

%==========================================================
\section{Preview of Course Themes}
\label{sec:course_preview}
%==========================================================

The remaining nine chapters of this textbook address the major contemporary
economic issues in the order determined by the pedagogical logic explained
in the preface. Each chapter builds on the institutional and analytical
foundations established here.

\textbf{Chapter~\ref{ch:inflation}: Inflation, Monetary Policy, and the Cost of Living.}
The immediate crisis of 2022--2023 examined in full: CPI measurement and its
biases, the aggregate demand--aggregate supply framework, the Phillips Curve,
the Bank of Canada's transmission mechanism, and the distributional consequences
of both inflation and the interest rate response.

\textbf{Chapter~\ref{ch:housing}: Canada's Housing Crisis.}
The economics of housing markets, the user cost model, the supply constraints
that prevent new construction from responding to demand, the role of
immigration and interest rates on the demand side, and the full range of
federal and provincial policy responses. This chapter flows directly from
Chapter~\ref{ch:inflation}: the interest rate hikes needed to control
inflation were the same rate hikes that deepened the mortgage crisis.

\textbf{Chapter~\ref{ch:labour}: Labour Markets, Employment, and Demographic Change.}
Labour supply and demand, labour force accounting, the Beveridge Curve, the
wage-productivity gap, minimum wage economics, non-standard employment, and
the relationship between an aging workforce and labour market tightness.

\textbf{Chapter~\ref{ch:inequality}: Income Inequality, Poverty, and Indigenous Economic Development.}
The Lorenz Curve and Gini coefficient, trends in Canadian income distribution
since 1980, intergenerational mobility, the gender wage gap, poverty
measurement, and the economic consequences of the historical dispossession
of Indigenous peoples.

\textbf{Chapter~\ref{ch:immigration}: Immigration, Population Growth, and the Canadian Economy.}
Canada's immigration system, the demographic arithmetic of the dependency
ratio, the labour market effects of immigration, the housing demand effects,
immigrant economic integration, and the policy debate over Canada's
historically high immigration targets.

\textbf{Chapter~\ref{ch:fiscal}: Fiscal Policy, Government Debt, and Public Finance.}
The federal budget process, deficit and debt arithmetic, the fiscal multiplier,
the COVID-era fiscal expansion, healthcare spending and its demographic drivers,
fiscal federalism and equalization, and the question of long-run fiscal
sustainability.

\textbf{Chapter~\ref{ch:trade}: International Trade, Supply Chains, and Canada-US Relations.}
Comparative advantage, the gains from trade, Canada's merchandise trade
patterns, CUSMA and its predecessors, the exchange rate and its effects,
the costs of tariff protection, supply chain vulnerabilities, and the
catastrophic potential of a US-Canada trade conflict.

\textbf{Chapter~\ref{ch:technology}: Technology, Automation, Artificial Intelligence, and Productivity.}
The Solow growth model, Canada's productivity puzzle, the task model of
automation, which Canadian jobs are at risk, the economics of AI adoption,
and what productivity growth requires in terms of investment, competition
policy, and innovation incentives.

\textbf{Chapter~\ref{ch:climate}: Energy, Climate Change, and Canada's Green Transition.}
The carbon externality as a market failure, the Pigouvian tax rationale
for carbon pricing, Canada's carbon price schedule, the oil sands and
their fiscal importance, the electricity grid transition, the distributional
incidence of carbon pricing, and the economic challenge of achieving
net-zero by 2050 while managing a responsible transition for fossil-fuel
dependent workers and communities.

%==========================================================
\section*{Chapter Summary}
%==========================================================
\addcontentsline{toc}{section}{Chapter Summary}

\begin{itemize}
  \item Contemporary economic issues are actively unfolding, contested, and
        consequential problems that require both theoretical tools and
        empirical evidence to analyse rigorously.

  \item Canada is the world's ninth-largest economy with a GDP of
        approximately \$3.0~trillion (2024), a population of 41.5~million,
        and a GDP per capita of roughly \$72,000~CAD (\$53,000~USD
        PPP-adjusted) --- placing it above the OECD average but
        significantly below the United States.

  \item GDP is measured using the expenditure approach: $Y = C + I + G +
        (X-M)$. Real GDP adjusts nominal GDP for changes in the price level,
        measured by the GDP deflator. Real GDP growth is the standard measure
        of economic expansion or contraction.

  \item Canada's economy is predominantly service-based (77\% of GDP), but
        with a disproportionately important natural resource sector (oil and
        gas, mining, forestry). Economically, Canada is dominated by Ontario
        and Quebec (60\% of national GDP), with Alberta as the third-largest
        provincial economy due to oil sands.

  \item Canada is highly trade-dependent (trade openness $\approx$ 60.5\%
        of GDP) and extremely concentrated in the US market (75\% of
        merchandise exports go to the United States).

  \item Canada's major economic institutions include the Bank of Canada
        (monetary policy, inflation targeting), the Department of Finance
        (fiscal policy), Statistics Canada (data), CMHC (housing), and
        provincial governments (healthcare, education, labour, natural
        resources).

  \item Monetary policy works through the overnight rate and its transmission
        channels (credit cost, asset prices, exchange rate, expectations).
        Fiscal policy works through government spending and taxation.
        The two interact: loose fiscal policy may require tight monetary
        policy, and vice versa.

  \item Economic reasoning requires: models (deliberate simplifications),
        positive--normative distinctions (what is vs.\ what should be),
        opportunity cost thinking, marginal analysis, and careful attention
        to the difference between correlation and causation.

  \item Canada's macroeconomic history from 2015 to 2026 --- summarized in
        Table~\ref{tab:macro_indicators} --- encompasses a commodity price
        shock (2015--16), a historic pandemic shock (2020), a rapid fiscal
        and monetary policy response, and an inflationary surge and
        tightening cycle (2022--2024) that will shape Canada's economy
        for years.
\end{itemize}

%==========================================================
\section*{Key Terms}
%==========================================================
\addcontentsline{toc}{section}{Key Terms}

\begin{description}[style=nextline, leftmargin=3cm]
  \item[\term{Aggregate demand}] Total spending on domestically produced
        goods and services in the economy ($C + I + G + NX$).
  \item[\term{Bank of Canada}] Canada's central bank; sets the overnight
        rate to achieve the 2\% inflation target.
  \item[\term{Budget deficit}] The excess of government spending over
        government revenues in a given fiscal year.
  \item[\term{Comparative advantage}] The ability to produce a good at
        a lower opportunity cost than another producer.
  \item[\term{Contemporary economic issues}] Actively unfolding economic
        problems that are contested and consequential for social well-being.
  \item[\term{Correlation}] A statistical relationship in which two
        variables tend to move together; does not imply causation.
  \item[\term{Fiscal policy}] The use of government spending and taxation
        to influence macroeconomic outcomes.
  \item[\term{GDP deflator}] The ratio of nominal GDP to real GDP multiplied
        by 100; a broad measure of economy-wide price level changes.
  \item[\term{Gross Domestic Product (GDP)}] The total market value of all
        final goods and services produced within a country's borders in a
        given period.
  \item[\term{GDP per capita}] GDP divided by population; a rough measure
        of average material living standards.
  \item[\term{Inflation targeting}] A monetary policy framework in which
        a central bank commits to keeping inflation within a specified range.
  \item[\term{Marginal cost and benefit}] The additional cost or benefit
        of one more unit of an activity.
  \item[\term{Model}] A deliberate simplification of reality used to
        focus analysis on the mechanisms that matter for a specific question.
  \item[\term{Monetary policy}] Actions by a central bank to influence
        interest rates and credit conditions, aiming at price stability
        and sustainable growth.
  \item[\term{Net exports (NX)}] Exports minus imports; the trade balance
        component of the GDP expenditure identity.
  \item[\term{Nominal GDP}] GDP measured in current-year prices.
  \item[\term{Opportunity cost}] The value of the next-best alternative
        foregone when making a choice.
  \item[\term{Overnight rate}] The interest rate targeted by the Bank of
        Canada; the rate at which major financial institutions lend to each
        other overnight.
  \item[\term{Pareto efficiency}] An allocation in which no one can be made
        better off without making someone else worse off.
  \item[\term{Positive economics}] Description and explanation of economic
        phenomena without value judgements; testable against evidence.
  \item[\term{Normative economics}] Value judgements about what economic
        outcomes \textit{should} be.
  \item[\term{Real GDP}] GDP measured in the prices of a fixed base year;
        removes the effect of inflation from output comparisons.
  \item[\term{Trade openness}] $(X + M)/Y \times 100$; the ratio of total
        trade to GDP.
  \item[\term{Transfer payments}] Government payments to individuals that
        do not correspond to the production of new goods or services
        (e.g.\ EI, CPP, OAS).
\end{description}

%==========================================================
\section*{Review Questions}
%==========================================================
\addcontentsline{toc}{section}{Review Questions}

\begin{enumerate}

\item \textbf{(Concept)} Define GDP and explain what it measures and what
it does not measure. Give two examples of economically important activities
that GDP omits, and explain why each omission might matter for policy.

\item \textbf{(Concept)} Explain the difference between nominal GDP growth
and real GDP growth. If Canada's nominal GDP rose from \$2.72~trillion
in 2022 to \$2.88~trillion in 2023, but the GDP deflator rose from
110.3 to 115.2, what was Canada's real GDP growth rate in 2023?
Show all steps.

\item \textbf{(Institutions)} Explain the difference between the Bank of
Canada and the Department of Finance. Why does Canada have two separate
institutions managing monetary and fiscal policy rather than one?
What problem does this separation solve?

\item \textbf{(Concept)} Classify each of the following statements as
positive or normative, and explain your classification:
\begin{enumerate}[(a)]
  \item ``A 25\% US tariff on Canadian auto exports would cost approximately
        80,000 Canadian jobs.''
  \item ``Canada should implement a national childcare program to increase
        female labour force participation.''
  \item ``Canada's GDP per capita is approximately 37\% lower than that
        of the United States.''
  \item ``The Bank of Canada should have started raising interest rates
        in late 2021 rather than waiting until March 2022.''
\end{enumerate}

\item \textbf{(Application)} Use Table~\ref{tab:macro_indicators} to
answer the following:
\begin{enumerate}[(a)]
  \item In which year was Canada's unemployment rate highest?
        What event caused this?
  \item In which year was Canada's inflation rate highest?
        What was the Bank of Canada's policy rate in that year?
  \item What happened to the federal debt-to-GDP ratio between 2019
        and 2020? Express the change in percentage points and explain
        what drove it.
\end{enumerate}

\item \textbf{(Critical Thinking)} Canada's trade openness ratio is
approximately 60.5\%, while the United States' is approximately 27\%.
Using the concept of opportunity cost, explain why a high trade openness
ratio creates both benefits and risks. Is there a ``right'' level of
trade openness for a country like Canada?

\end{enumerate}

%==========================================================
\section*{Quantitative Problems}
%==========================================================
\addcontentsline{toc}{section}{Quantitative Problems}

\begin{enumerate}

\item \textbf{GDP Calculation.} You are given the following data for a
hypothetical Canadian province:
\begin{center}
\begin{tabular}{lr}
\toprule
Item & Value (\$B) \\
\midrule
Household consumption expenditure & 480 \\
Business investment in fixed capital & 120 \\
Change in business inventories & 8 \\
Provincial government spending on goods/services & 140 \\
Exports to other provinces and countries & 210 \\
Imports from other provinces and countries & 195 \\
\bottomrule
\end{tabular}
\end{center}
\begin{enumerate}[(a)]
  \item Calculate nominal GDP using the expenditure approach
        (Equation~\ref{eq:gdp_expenditure}). Show all steps.
  \item What is the trade balance for this province?
        What does its sign indicate?
  \item If the population is 4.8~million, calculate GDP per capita.
\end{enumerate}

\item \textbf{Real vs. Nominal GDP.} The following data describe a
simplified economy producing only two goods: wheat and lumber.
\begin{center}
\begin{tabular}{lcccc}
\toprule
& \multicolumn{2}{c}{\textbf{Year 1 (Base)}} & \multicolumn{2}{c}{\textbf{Year 2}} \\
\cmidrule(lr){2-3}\cmidrule(lr){4-5}
Good & Price & Quantity & Price & Quantity \\
\midrule
Wheat  & \$4.00/bushel & 500 & \$5.50/bushel & 520 \\
Lumber & \$8.00/board-ft & 300 & \$9.00/board-ft & 330 \\
\bottomrule
\end{tabular}
\end{center}
\begin{enumerate}[(a)]
  \item Calculate nominal GDP in Year 1 and Year 2.
  \item Calculate real GDP in Year 2 using Year 1 prices.
  \item Calculate the nominal GDP growth rate and the real GDP growth rate
        between Year 1 and Year 2.
  \item What is the GDP deflator in Year 2 (with Year 1 = 100)?
        What is the implied inflation rate between Year 1 and Year 2?
  \item If the nominal growth rate was 17.8\% but the real growth rate
        was only 6.5\%, what does the difference represent?
\end{enumerate}

\item \textbf{Trade Openness.} Using Table~\ref{tab:trading_partners}
and the following data:
\begin{itemize}
  \item Canada's total merchandise exports, 2024: \$897~billion CAD
  \item Canada's total merchandise imports, 2024: \$918~billion CAD
  \item Canada's nominal GDP, 2024: \$3,001~billion CAD
\end{itemize}
\begin{enumerate}[(a)]
  \item Calculate Canada's trade openness ratio
        (Equation~\ref{eq:trade_openness}).
  \item Calculate Canada's merchandise trade balance.
  \item What share of Canada's total exports goes to the United States?
        (Use Table~\ref{tab:trading_partners}.)
  \item If the United States were to impose a 25\% tariff on all
        Canadian goods, and this caused Canada's US-bound exports to
        fall by 30\%, what would be the new trade balance?
        (Assume all other trade is unchanged.)
\end{enumerate}

\item \textbf{GDP per Capita Growth.} The relationship between GDP per
capita growth, GDP growth, and population growth is approximately:
\[
  g_{\text{per capita}} \approx g_{\text{GDP}} - g_{\text{population}}
\]
\begin{enumerate}[(a)]
  \item In 2023, Canada's real GDP grew by 1.2\% and its population
        grew by approximately 3.2\% (driven largely by immigration).
        What was the approximate growth rate of real GDP \textit{per capita}?
        What does this imply about average living standards?
  \item If Canada's potential GDP growth rate is 2.0\% and population
        growth is 3.0\%, what is the implied trend in GDP per capita?
        Is this sustainable?
  \item Suggest two economic policies that could raise Canada's GDP
        growth rate so that it exceeds its population growth rate.
        For each policy, identify a potential trade-off.
\end{enumerate}

\item \textbf{Opportunity Cost.} The federal government has \$20~billion
in additional fiscal capacity. Three competing proposals are on the table:
\begin{itemize}
  \item Proposal A: Invest in new public housing construction.
  \item Proposal B: Cut the corporate income tax rate by 2 percentage points.
  \item Proposal C: Fund a national pharmacare program for uninsured drugs.
\end{itemize}
\begin{enumerate}[(a)]
  \item For each proposal, identify one major benefit and one major cost
        (the opportunity cost of choosing that proposal over the others).
  \item Which economic concept (efficiency, equity, or both) is most
        relevant for evaluating each proposal?
  \item Is the choice among these three proposals a positive economic
        question, a normative question, or both? Explain.
\end{enumerate}

\item \textbf{Data Exercise.} Visit the Statistics Canada website
(\url{www.statcan.gc.ca}) and locate Table 36-10-0104-01 (GDP by
expenditure, seasonally adjusted, at annual rates).
\begin{enumerate}[(a)]
  \item Report the real GDP growth rate (annualized) for the most recent
        available quarter.
  \item Break down GDP by component (C, I, G, NX) for the most recent quarter.
        Which component is largest? Which is most volatile over the past
        four quarters?
  \item Calculate the percentage change in each component over the past
        four quarters.
  \item Based on the current composition of growth, assess whether Canada's
        economic expansion is likely to be sustainable. Which components
        would you expect to change if interest rates rise further?
\end{enumerate}

\end{enumerate}

%==========================================================
\section*{Essay Questions}
%==========================================================
\addcontentsline{toc}{section}{Essay Questions}

\begin{enumerate}

\item \textbf{(500--700 words)} Using data from Table~\ref{tab:macro_indicators},
write an essay describing the macroeconomic history of Canada between 2019
and 2026. Your essay should identify the major shocks, explain the policy
responses, assess their effectiveness, and discuss the economic legacy of
this period for Canada going forward. Use at least three specific data
points from the table to support your argument.

\item \textbf{(400--600 words)} Canada runs an enormous merchandise trade
surplus with the United States but trade deficits with most other countries
(see Table~\ref{tab:trading_partners}). Using the concepts of comparative
advantage and opportunity cost, explain what this pattern tells us about
Canada's productive specialization. Discuss one risk and one benefit of
Canada's trade concentration in the US market.

\item \textbf{(400--500 words)} The opening vignette of this chapter
describes four people --- Diane, Marcus, the sawmill manager, and Priya
--- each facing an economic challenge shaped by macroeconomic forces.
Choose two of these four people. For each, identify the specific macroeconomic
force most relevant to their situation, explain the economic mechanism
at work, and identify one policy from Table~\ref{tab:institutions}
(or one of the institutions described in Section~\ref{sec:institutions})
that is most directly relevant to addressing their challenge.

\end{enumerate}

%==========================================================
\section*{Discussion Questions}
%==========================================================
\addcontentsline{toc}{section}{Discussion Questions}

\begin{enumerate}

\item Figure~\ref{fig:gdp_per_capita} shows that Canada's GDP per capita
is approximately 37\% lower than that of the United States, despite
Canada having a comparable education system, similar institutions, and
vast natural resource wealth. What do you think explains this gap?
Is it a problem that Canada should try to close? What would closing
it require?

\item Should the Bank of Canada be fully independent from the elected
federal government? What would you gain, and what would you lose,
by giving Parliament more direct control over monetary policy?
Draw on the Policy Debate box in Section~\ref{sec:policy_toolkit}
in your answer.

\item Canada's regional economic inequality --- illustrated in
Figure~\ref{fig:gdp_by_province} --- means that national economic
policies have unequal regional consequences. Is this a problem
unique to Canada, or a challenge faced by all federations?
How should the federal government balance national economic objectives
with regional equity?

\item The opening of this chapter describes four people whose lives
are shaped by macroeconomic forces they cannot individually control.
Do you think most Canadians understand how economics affects their
daily lives? What are the consequences --- political, social, and
economic --- of widespread economic illiteracy?

\end{enumerate}
